\begin{claim}{D.2.36}{Proved}{Every entire fourth-suspension unit has zero linear image}
Put $B=K(\FF_2,2)$, $X=\Sigma B$, $A=QX$, and
$Y=\Sigma^2B$. Every ordinary pointed map
\[
             h:\Sigma^5B\longrightarrow\Sigma^2B
\]
induces the zero map on degree-eight mod-two homology.
The same is therefore true of every ordinary pointed map
$h:\Sigma^4A\to Y$.

For the double adjunct $G:\Sigma^2A\to T=\Omega^2Y$,
write $B_0=\sigma b\in H_4A$, where $b$ is dual to
$v=\operatorname{Sq}^1u\in H^3B$. Then
\[
                   G_*(\sigma^2B_0)=\epsilon k,
                   \qquad \epsilon\in\FF_2,
                   \quad k=\{a,b\}.
                   \tag{D.2.36a}
\]
In particular, the alternative bottom pattern
$b^2+k+p_3$ cannot occur on an entire unit. The pure
commutator pattern is realized on the entire unit by D.2.29;
its extension over the full infinite-loop source is not
settled by this statement.
\end{claim}
\begin{proof}
All coefficients are $\FF_2$. Remove the suspensions from
the induced cohomology homomorphism to obtain a linear map
\[
 \phi=\sigma^{-5}h^*\sigma^2:
 \widetilde H^dB\longrightarrow\widetilde H^{d-3}B.
                   \tag{D.2.36b}
\]
It commutes with every Steenrod square, by naturality and
suspension compatibility. It is not assumed multiplicative.
Serre's theorem gives
\[
 H^*B=\FF_2[u_2,v_3,w_5,t_9,s_{17},\ldots],\qquad
 v=\operatorname{Sq}^1u,\quad w=\operatorname{Sq}^2v,
 \quad t=\operatorname{Sq}^4w,\quad s=\operatorname{Sq}^8t;
                   \tag{D.2.36c}
\]
the later generator degrees are $2^j+1$. The exact source
and operation conventions are
\cite[\S1, nos.~2.1--2.4; \S2, no.~8, Theorem~2,
and no.~9, pp.~199, 203--204]{I_Serre1953}.

Suppose $\phi(u^3)=v$. The following five identities follow
directly from Cartan and instability:
\[
\begin{aligned}
 \operatorname{Sq}^4(u^2)&=u^4,&
 \operatorname{Sq}^3(uv)&=uv^2+vw,\\
 \operatorname{Sq}^2(u^3)&=u^4+uv^2,&
 \operatorname{Sq}^6(v^2)&=v^4,\\
 \operatorname{Sq}^4(vw)&=vt+v^4.
\end{aligned}
                   \tag{D.2.36d}
\]
Since $\phi(u^2)\in H^1B=0$, the first gives
$\phi(u^4)=0$. The third now gives $\phi(uv^2)=w$.
The second gives $\phi(vw)=w$, since
$\operatorname{Sq}^3\phi(uv)=0$ by instability in degree two.
Likewise $\operatorname{Sq}^6\phi(v^2)=0$ in degree three,
so the last two identities give
\[
                       \phi(v^4)=0,\qquad \phi(vt)=t.
                       \tag{D.2.36e}
\]

Four more explicit source identities finish the contradiction:
\[
\begin{aligned}
 \operatorname{Sq}^{10}(uvw)&=u^2v^2w^2,&
 \operatorname{Sq}^9(ut)&=ut^2+vs,\\
 \operatorname{Sq}^8(vt)&=vs,&
 \operatorname{Sq}^4(uwt)&=ut^2+u^2v^2w^2.
\end{aligned}
                       \tag{D.2.36f}
\]
For clarity, the generator formulas used in Cartan here are
\[
 \operatorname{Sq}^1t=w^2,\quad
 \operatorname{Sq}^8t=s,\quad
 \operatorname{Sq}^9t=t^2,\quad
 \operatorname{Sq}^jt=0\quad(2\leq j\leq7).
                       \tag{D.2.36g}
\]
The first and last follow by Adem applied to $t=\operatorname{Sq}^4w$,
using $\operatorname{Sq}^1w=v^2$ and
$\operatorname{Sq}^2w=\operatorname{Sq}^3w=0$;
the middle two are the generator definition and the top-square
identity. Thus every term in \textup{(D.2.36f)} is specified.

The first identity of \textup{(D.2.36f)} has zero image under
$\phi$, since $\phi(uvw)$ has degree seven. The second has
zero image as a ninth square of a degree-eight class.
The third, together with \textup{(D.2.36e)}, gives
\[
 \phi(u^2v^2w^2)=0,\qquad
 \phi(ut^2)=\phi(vs)=s.
                       \tag{D.2.36h}
\]
The final identity would therefore give
\[
                       \operatorname{Sq}^4\phi(uwt)=s.
                       \tag{D.2.36i}
\]
But $\phi(uwt)$ has degree thirteen, where every class of
$H^*B$ is decomposable: there is no algebra generator in
that degree. Cartan preserves the ideal of decomposables,
whereas $s$ is an indecomposable generator. This contradiction
proves $\phi(u^3)=0$. Also
\[
 \phi(v^2)=\phi(\operatorname{Sq}^2u^2)
          =\operatorname{Sq}^2\phi(u^2)=0.
\]
The degree-six basis is $u^3,v^2$, so duality proves the
degree-eight homology assertion for the unit map.

The unit $X\to QX$ is an isomorphism on degree-four
homology: the first extra stable weight starts in degree six.
After four suspensions it therefore spans all degree-eight
source homology. Restriction to this actual unit proves the
same vanishing for every $\Sigma^4A\to Y$.

Finally, the complete primitive basis of $H_6T$ is
$b^2,k,p_3$, as in D.2.24. Double homology suspension sends
$p_3$ to the linear $f$ label and kills the other two.
The homology vanishing just proved excludes its coefficient.
Since $\operatorname{Sq}_*^2(\sigma^2B_0)=0$, while
$\operatorname{Sq}_*^2b^2=a^2$ and
$\operatorname{Sq}_*^2k=0$, the $b^2$ coefficient is zero
as well. This proves \textup{(D.2.36a)} without imposing
any multiplicativity on the ordinary map.

The exact program \texttt{code/D2\_k4\_linear\_module.py}
independently constructs all source monomials, uses Adem and
Cartan, and allows every coefficient of $\phi$ except the
single prescribed value $\phi(u^3)=v$. Its degree-twenty
inconsistency certificate is precisely the nine identities
above. The range through degree nineteen is consistent as
a truncated module problem; no realization of that truncated
solution is asserted.
\end{proof}
